\documentclass[10pt,a4paper]{amsart}
\usepackage[margin=19mm]{geometry}
\usepackage{amsmath,amssymb,mathtools}
\usepackage[scr=rsfs,cal=euler]{mathalfa}
\usepackage{xfrac}
\usepackage[unicode,hidelinks,bookmarks=false]{hyperref}
\newcommand{\ie}{i.e.\ }
\newcommand{\cf}{cf.\ }
\newcommand{\resp}{respectively\ }
\newcommand{\Subsection}{\subsection}
\providecommand{\nobreakdash}{\nobreak\hbox{-}}
\setlength{\emergencystretch}{2em}
\multlinegap0pt
\usepackage{xfrac}
\usepackage{mathtools}
\usepackage{amssymb}
\usepackage{tikz}
\newtheorem{lem}{Lemma}
\DeclareMathOperator{\An}{\text{An}\Omega}
\newcommand{\quand}{\quad\text{and}\quad}
\newcommand{\spacetime}{(M,\Omega,g,\nabla)}
\title{On Twisted Spacetimes: a new class of Galilean cosmological models}
\author{Daniel de la Fuente, Rafael M. Rubio, Jose Torrente-Teruel}
\date{Source: arXiv:2406.10155v1 (CC BY 4.0)}
\begin{document}
\maketitle

\noindent\emph{Source and licence.} On Twisted Spacetimes: a new class of Galilean cosmological models. Version arXiv:2406.10155v1 (CC BY 4.0), distributed by arXiv under the Creative Commons Attribution 4.0 International licence, http://creativecommons.org/licenses/by/4.0/, as recorded in the arXiv licence metadata for this exact version and shown on the abstract page https://arxiv.org/abs/2406.10155v1. Full licence notice: \texttt{licenses/arxiv-2406.10155-CC-BY-4.0.txt}.\par\smallskip

\noindent\textit{Editorial scope.} Counted unit: the article's lem eq:alpha (source0.tex lines 714--723). The article's complete original proof of it is delivered with it (source0.tex lines 724--739, one \texttt{proof} environment of 16 lines). No mathematics is written, repaired or re-derived: the statement and the proof are the article's own lines, in the article's own notation, and no retained line is substituted or re-flowed.  Retained uncounted as the source-local context the proof needs: lines 707--709.  Nothing was omitted from any retained span, so the package carries no omission record.  No retained line contains a citation, so the wrapper carries no bibliography and no bibliography entry.  No retained line contains a figure environment, an image-inclusion command or drawing code, so no image is delivered and the package declares no asset dependency.  Editorial formatting only, in addition: an AMS \texttt{amsart} preamble that restates the article's own declarations and macros used by the retained lines, namely the environments \texttt{lem} on separate counters, together with the standard packages those lines need.  The retained environments print in this document as Lemma 1 in order of appearance, whereas the article prints the same items with the numbering of its own full document.  The complete native source of the article as distributed in the arXiv e-print for this exact version is bundled unchanged as \texttt{sources/160824/source0.tex}.
    \begin{equation}\label{eq:torqued}
        \nabla_X K = \rho X + \alpha(X) K \quand \alpha(K) = 0, \quad \forall X \in \Gamma(T M)\,,
    \end{equation}

\begin{lem}
    Let $\spacetime$ be a spacetime and suppose it admits a Galilean torqued vector field $K\in \Gamma(TM)$ with torqued function $\rho\in C^\infty(M)$ and torqued form $\alpha \in \bigwedge^1(M)$. Then, 
    \begin{equation}\label{eq:alpha}
        \alpha(V) = \tfrac{1}{\Omega(K)} V\big(\Omega(K)\big), \quad\forall V\in \Gamma(\An)\,.
    \end{equation}
 and
    \begin{equation}\label{eq:rho}
        \rho = \tfrac{1}{\Omega(K)} K\big( \Omega(K)\big)\,.
    \end{equation}
\end{lem}

\begin{proof}
    Consider the field of observers $Z=\frac{1}{\Omega(K)} K$. Then, for any $V\in \Gamma(\An)$,
    \begin{equation}\label{eq:proof derivative Z}
    \begin{split}
        \nabla_V Z ={}& V\big(\tfrac{1}{\Omega(K)}\big) K + \tfrac{1}{\Omega(K)} \nabla_V K \\
        ={}& \tfrac{\rho}{\Omega(K)} V + \left[ \tfrac{\alpha(V)}{\Omega(K)} + V\big(\tfrac{1}{\Omega(K)}\big) \right] K\,.
    \end{split}
    \end{equation}
    Now, making use of the compatibility of the connection with $\Omega$ we have $\Omega(\nabla_V Z) = V(\Omega(Z)) = 0$. Then, from (\ref{eq:proof derivative Z}) we get (\ref{eq:alpha}).

    On the other hand, from equation (\ref{eq:torqued}) we have $\nabla_K K = \rho K$. Consequently,
    \[
        K(\Omega(K)) = \Omega(\nabla_K K) = \rho \Omega(K)\,.
    \]
    Thus, it follows expression (\ref{eq:rho}).
\end{proof}

\end{document}
